\subsection{Irreducible Torsion and Higher Winding Sectors}
  \label{sec:irreducible-torsion}

  Within the spectral framework established in Chapter~\ref{sec:spectral-admissibility},
  the condition
  \(
  [\Delta^{(0)}_G,\Omega_w]\neq 0
  \)
  signals spectral frustration for all $w\ge 2$.
  However, non-commutativity alone does not distinguish the qualitative
  structure of different higher-winding sectors.
  A further structural criterion is required.

  \subsubsection*{From Anisotropy to Irreducibility}

    In the $w=2$ sector, torsion induces anisotropy within the projection
    fiber, leading to a controlled splitting
    \(
    \Pi = \Pi_{\parallel}\oplus\Pi_{\perp}.
    \)
    The associated dynamical selection principle therefore operates
    within a reducible decomposition of the first spectral multiplet.

    For $w\ge 3$, the situation may be qualitatively different.
    The torsional operator $\Omega_w$ can act on the admissible subspace
    in a manner that is not decomposable into independent lower-winding
    sectors.

    Formally, let
    \begin{equation}
      \mathcal{A}_w
      \equiv
      \mathrm{alg}\!\left(
                      \Delta^{(0)}_G,\,\Omega_w
      \right)
      \Big|_{H_{\mathrm{adm}}^{(w)}}
    \end{equation}
    denote the algebra generated by relaxation and torsion,
    restricted to the admissible configuration space.

    If the commutant
    \begin{equation}
      \mathcal{A}_w'
      =
      \big\{
      X\in\mathrm{End}(H_{\mathrm{adm}}^{(w)})
      \mid
      [X,A]=0\ \forall A\in\mathcal{A}_w
      \big\}
    \end{equation}
    reduces to scalar multiples of the identity,
    the torsional action is irreducible in the admissible sector.

    Irreducibility implies that no non-trivial projector exists that
    simultaneously preserves admissibility and diagonalizes the torsional
    action.
    In that case a naive decomposition
    \(
    \Omega_w \simeq \Omega_{1}^{\oplus w}
    \)
    is not dynamically accessible within $H_{\mathrm{adm}}^{(w)}$.

    Operationally, irreducibility can be tested by computing the
    dimension of the commutant $\mathcal{A}_w'$.
    If $\dim \mathcal{A}_w' = 1$, the torsional sector is algebraically
    irreducible within $H_{\mathrm{adm}}^{(w)}$.

  \subsubsection*{Admissibility as a Structural Constraint}

    The admissible subspace $H_{\mathrm{adm}}^{(w)}$ is defined by the
    bounded-relaxation condition.
    Factorizations that would decompose $\Omega_w$ into independent
    lower-winding components may require intermediate configurations
    that violate admissibility.

    In such cases irreducibility is not an additional axiom but a
    consequence of bounded relaxation:
    the decomposition path exits the admissible domain and is therefore
    dynamically forbidden.

    This mechanism provides a structural explanation for why higher
    winding sectors may exhibit stronger mass amplification than the
    $w=2$ case.
    The obstruction is therefore not merely energetic but algebraic.

  \subsubsection*{Spectral Consequences for $w=3$}

    The $w=3$ sector is the first case in which genuine algebraic
    irreducibility may arise.
    If $\mathcal{A}_3'$ is trivial within $H_{\mathrm{adm}}^{(3)}$,
    the torsional action cannot be reduced to a sum of independent
    leptonic sectors.

    In that regime the relevant mass scale is controlled by a global
    spectral invariant rather than by a two-mode dynamical selection rule.
    A natural candidate is the torsional action
    \(
    \mathcal{A}(3)
    \)
    defined in Eq.~\eqref{eq:Aw-fredholm}, or an associated canonical
    ratio extracted from the spectrum of
    \(
    \Delta^{(0)}_G+\Omega_3
    \)
    restricted to $H_{\mathrm{adm}}^{(3)}$.

    Identifying such an invariant constitutes the next structural step
    beyond the leptonic sector.
